\chapter{Representation Learning}\label{ch:ml:representation-learning}

A desk has years of order-book data for an instrument and labels for a few weeks of it: the weeks since it began trading
through a new venue, the only period the label it cares about exists for. The unlabelled years are not useless. They
show how books move, which states follow which, what is noise and what persists, and a model can learn that before it
sees a single label. On the chapter's simulated sessions, an encoder pre-trained to forecast the book's own next few
seconds, then read by a \omterm{def:ml:representation-learning:probe}{linear probe}, beats every model trained on the raw book from the labels alone when there are
300 or 1\,000 labelled windows; an encoder pre-trained to reconstruct its input does not. This chapter builds the
self-supervised objectives, measures which of them help and when, and explains why reconstruction is the wrong pretext
for a forecast.

\section{Representations and autoencoders}

\begin{definition}[Representation learning]\label{def:ml:representation-learning:representation}
\emph{Representation learning}\index{representation learning} fits a map $h = e_\theta(x)$ from raw inputs to a
lower-dimensional code $h$ on which later tasks are easier; the map is judged by what those tasks achieve with it.
\end{definition}

Chapter~9 defined the \omterm{def:ml:cross-sectional-deep-models:ae}{autoencoder}: an encoder, a decoder, and a reconstruction loss. Its oldest lesson is also its most
useful one here.

\begin{proposition}[A linear autoencoder is PCA]\label{prop:ml:represent:pca}
Among linear encoders $h = Ax$ and decoders $\hat x = Bh$ with a $k$-dimensional code, the reconstruction error
$\sum_i\lVert x_i - BAx_i\rVert^2$ is minimised when $BA$ is the orthogonal projection onto the span of the first $k$
principal components of the (centred) data (Baldi and Hornik, 1989).
\end{proposition}

\begin{proof}
$BAX$ has rank at most $k$, and by the Eckart--Young theorem the best rank-$k$ approximation of $X$ in Frobenius norm is
its truncated singular value decomposition, the projection onto the top $k$ right singular vectors, which are the
principal directions (Book~4, chapters~22 and~25). That projection is of the form $BA$.
\end{proof}

A reconstruction objective keeps the directions of largest variance, whatever their relation to the task. On an order
book the largest variation is in the sizes of the deep levels, which say little about the next seconds; the
imbalance at the top of the book, which says a great deal, is a small part of the variance.

\begin{definition}[Denoising autoencoder]\label{def:ml:representation-learning:dae}
A \emph{denoising autoencoder}\index{denoising autoencoder} reconstructs the clean input from a corrupted copy (entries
masked to zero, Gaussian noise added), so that the code must capture the input's structure rather than copy it (Vincent
et al., 2008).
\end{definition}

\section{Self-supervised objectives}

\begin{definition}[Self-supervised learning, pre-training, fine-tuning]\label{def:ml:representation-learning:ssl}
\emph{Self-supervised learning}\index{self-supervised learning} trains on a pretext task whose targets are made from the
unlabelled data themselves: a masked part of the input, the input's future, or another view of the same example.
\emph{Pre-training}\index{pre-training} is such a fit, done before the task's labels are used; \emph{fine-tuning}\index{fine-tuning}
continues training the pre-trained network, with a new output head, on the labelled task, usually at a small \omterm{def:ml:trees-and-boosting:boosting}{learning
rate}.
\end{definition}

\begin{definition}[Linear probe]\label{def:ml:representation-learning:probe}
A \emph{linear probe}\index{linear probe} is a linear model (here a ridge regression) fitted on the frozen codes of a
pre-trained encoder; its score measures how much of the task the representation makes linearly available (Alain and
Bengio, 2017).
\end{definition}

Forecasting the input's own future is the pretext that language models use, and it is available on any market data:
the book a few seconds later is not a label, it is the next part of the same stream. The chapter's third encoder is
pre-trained to predict the change of the whole ten-level snapshot two seconds after the window ends. The label the desk
cares about (the mean mid over the next ten seconds) is never used in \omterm{def:ml:representation-learning:ssl}{pre-training}.

\section{Contrastive learning and augmentations for market data}

\begin{definition}[Contrastive learning, data augmentation]\label{def:ml:representation-learning:contrastive}
\emph{Data augmentation}\index{data augmentation} makes new training inputs by transformations that should not change
what matters about them (noise, masking, small shifts). \emph{Contrastive learning}\index{contrastive learning} trains an
encoder so that two augmented views of the same example have similar codes and views of different examples dissimilar
ones; with the InfoNCE loss each view must pick its partner among the $2m - 1$ other views of a batch of $m$ examples,
through a softmax of cosine similarities divided by a temperature (van den Oord et al., 2018; Chen et al., 2020).
\end{definition}

The choice of augmentations is the model's prior about what does not matter. Image models use crops and colour changes;
for an order book there is no agreed set. The chapter's augmentations add Gaussian noise of 0.3 standard deviations and
mask 20\% of the entries of a window. Both are plausible, and both turn out to preserve the wrong things: the codes they
produce are no better for the forecast than the \omterm{def:ml:cross-sectional-deep-models:ae}{autoencoder}'s.

\section{Pre-training, fine-tuning and the linear probe}

The chapter's data are ten-minute \texttt{firm.tape} sessions: sixteen unlabelled ones for \omterm{def:ml:representation-learning:ssl}{pre-training}, four labelled
ones, one for \omterm{def:ml:trees-and-boosting:early}{early stopping} and six for testing. An input is a window of ten half-second snapshots of ten levels (400
numbers); the target is chapter~8's forward change (the mean mid over the next ten seconds against the current mid),
here as a number, scored by its information coefficient on each test session. The three encoders (an MLP with 64
hidden units and a 16-dimensional code) are pre-trained for ten \omterm{def:ml:neural-networks-for-noisy-tabular-data:backprop}{epochs}; each labelled-set size is averaged over four
draws, the first windows of each labelled session.

\begin{omfigure}
\begin{tikzpicture}
\begin{semilogxaxis}[width=11cm,height=6cm,xlabel={labelled windows},ylabel={test information coefficient},
  tick label style={font=\small},label style={font=\small},xmin=220,xmax=6000,ymin=0,ymax=0.6,
  xtick={300,1000,4444},xticklabels={300,1\,000,4\,444},grid=major,grid style={gray!25},legend columns=3,
  legend cell align=left,legend style={font=\footnotesize,at={(0.5,-0.24)},anchor=north,draw=none,fill=none,
  /tikz/every even column/.append style={column sep=6pt}},table/col sep=comma]
\addplot[gray,very thick,mark=triangle*] table[x=windows,y=raw_ridge]{figdata/ml/10-representation-learning/labels.csv};
\addplot[omProp,very thick,mark=diamond*] table[x=windows,y=from_scratch]{figdata/ml/10-representation-learning/labels.csv};
\addplot[omThm!70,thick,mark=square] table[x=windows,y=autoencoder_probe]{figdata/ml/10-representation-learning/labels.csv};
\addplot[omThm,thick,dashed,mark=o] table[x=windows,y=contrastive_probe]{figdata/ml/10-representation-learning/labels.csv};
\addplot[omDef,very thick,mark=*] table[x=windows,y=predictive_probe]{figdata/ml/10-representation-learning/labels.csv};
\addplot[omDef,thick,dashed,mark=*] table[x=windows,y=predictive_fine_tuned]{figdata/ml/10-representation-learning/labels.csv};
\addplot[black!70,thick,dotted,mark=x] table[x=windows,y=features_ridge]{figdata/ml/10-representation-learning/labels.csv};
\legend{ridge on the raw window,network from scratch,autoencoder probe,contrastive probe,predictive probe,
predictive fine-tuned,ridge on hand-made features}
\end{semilogxaxis}
\end{tikzpicture}

\omcaption{Test IC of the forecast of the next ten seconds against the number of labelled windows (300 and 1\,000: averages
over four draws; 4\,444: all four labelled sessions). The pre-trained encoders never saw a label. Data:
\texttt{ml\_represent.compare}.}\label{fig:ml:represent:labels}
\end{omfigure}

\Cref{fig:ml:represent:labels} has one winner in the low-label regime and one lesson throughout. With 300 labelled
windows (two and a half minutes of trading), the probe on the predictive encoder scores 0.22, against 0.12 for ridge on
the raw window and 0.18 for the same network trained from scratch; with 1\,000 windows, 0.29 against 0.24 and 0.20. With
all 4\,444 labelled windows, ridge on the raw window has caught up (0.45 against the probe's 0.40), and only \omterm{def:ml:representation-learning:ssl}{fine-tuning}
the predictive encoder keeps level with it (0.46), still ahead of the network from scratch (0.38). The reconstruction and
contrastive codes do worse than the raw window at 1\,000 and 4\,444 labels, and at 300 the \omterm{def:ml:cross-sectional-deep-models:ae}{autoencoder} is barely better
(0.15) and the contrastive encoder worse (0.05), as \cref{prop:ml:represent:pca} warns: they keep what varies, not what
predicts.

The hand-made order-flow features of chapter~8 beat everything once there are 1\,000 labels (0.50 and 0.55), and are the
worst of all with 300 (0.05): ten features and a ridge regression fitted on two and a half minutes of flow that rarely
moved. A label-free \omterm{def:ml:representation-learning:ssl}{pre-training} teaches a network what microstructure research has already written down; its value is
largest where that knowledge is missing (a new market, a new data type) and the labels are few.

\begin{method}[Pre-training on market data]\label{met:ml:represent:method}
\begin{enumerate}
\item Choose a pretext close to the task: forecasting the stream's own future beats reconstructing it for a forecast.
\item Pre-train on sessions disjoint from the labelled and test ones; fit normalisation on the \omterm{def:ml:representation-learning:ssl}{pre-training} data.
\item Measure with a \omterm{def:ml:representation-learning:probe}{linear probe} first; fine-tune only if the probe is useful, at a small \omterm{def:ml:trees-and-boosting:boosting}{learning rate}, with \omterm{def:ml:trees-and-boosting:early}{early
stopping}.
\item Report the learning curve in labelled data against ridge on the raw input and a network from scratch: \omterm{def:ml:representation-learning:ssl}{pre-training}
is worth its cost only where the curves separate.
\end{enumerate}
\end{method}

\section{Tutorial: learning without labels}\label{tut:ml:representation-learning:tut}

\begin{tutorial}
\textbf{Goal.} Pre-train three encoders on unlabelled sessions and measure their codes with probes and \omterm{def:ml:representation-learning:ssl}{fine-tuning} as the
number of labels grows. \textbf{End state:} \cref{fig:ml:represent:labels}.
\begin{tutsteps}
\item \textbf{The contrastive objective}: two augmented views, InfoNCE over the batch.
\omcode{code/firm/represent/firm_represent.py}{71}{97}{InfoNCE and contrastive pre-training.}\label{lst:ml:represent:contrastive}
\item \textbf{Self-supervised forecasting}: the encoder predicts a label-free future of its input.
\omcode{code/firm/represent/firm_represent.py}{100}{115}{Pre-training by forecasting the stream's own future.}\label{lst:ml:represent:predictive}
\item \textbf{Run} \texttt{ml\_represent.compare(300)}, \texttt{compare(1000)}, \texttt{compare(0)} and
\texttt{fig\_represent.py}.
\end{tutsteps}
\textbf{What to change next.} Replace the contrastive augmentations by a time shift of one snapshot (two nearby windows as a
positive pair) and see whether the codes improve; forecast the snapshot eight seconds ahead instead of two.
\end{tutorial}

\section{Build: representation learning}\label{bld:ml:representation-learning:represent}

\begin{build}
\bfield{Purpose} Encoders pre-trained on the firm's unlabelled market data, measured before any model is built on them.
\bfield{Interface} \texttt{Encoder(d\_in, code, hidden)}, \texttt{augment(x, gen, noise, mask)},
\texttt{pretrain\_dae}, \texttt{pretrain\_contrastive}, \texttt{pretrain\_predictive(X, Y, code, hidden, epochs, seed)},
\texttt{info\_nce(z1, z2, tau)}, \texttt{encode(enc, X)}, \texttt{linear\_probe(Z, y, Zt, alpha)},
\texttt{fine\_tune(enc, X, y, Xv, yv, lr, epochs, patience, seed)}, \texttt{scratch(X, y, Xv, yv, code, hidden)}.
\bfield{Rules} \omterm{def:ml:representation-learning:ssl}{Pre-training} data disjoint from labelled and test data; every random draw from the run's seed; one CPU
thread.
\bfield{Acceptance tests} \texttt{code/firm/represent/tests/}: the InfoNCE loss is $\ln(2m-1)$ for uninformative codes and
near zero for identical views; a linear \omterm{def:ml:representation-learning:dae}{denoising autoencoder}'s code spans the top principal subspace; the predictive
encoder's probe recovers a planted relation between a window and its future; \omterm{def:ml:representation-learning:ssl}{fine-tuning} does not modify the pre-trained
encoder passed in.
\bfield{Stretch} Masked-snapshot \omterm{def:ml:representation-learning:ssl}{pre-training} with a small \omterm{def:ml:sequence-models-on-order-books:attention}{transformer}; codes pooled across several simulated
instruments.
\end{build}

\begin{omsources}
\item G.~E.~Hinton and R.~R.~Salakhutdinov, ``Reducing the dimensionality of data with neural networks'', \emph{Science}
313, 2006.
\item P.~Vincent, H.~Larochelle, Y.~Bengio and P.-A.~Manzagol, ``Extracting and composing robust features with denoising
autoencoders'', \emph{ICML}, 2008.
\item P.~Baldi and K.~Hornik, ``Neural networks and principal component analysis: learning from examples without local
minima'', \emph{Neural Networks} 2(1), 1989.
\item A.~van den Oord, Y.~Li and O.~Vinyals, ``Representation learning with contrastive predictive coding'', arXiv, 2018;
T.~Chen, S.~Kornblith, M.~Norouzi and G.~Hinton, ``A simple framework for contrastive learning of visual representations'',
\emph{ICML}, 2020.
\item Z.~Yue et al., ``TS2Vec: towards universal representation of time series'', \emph{AAAI}, 2022.
\item G.~Alain and Y.~Bengio, ``Understanding intermediate layers using linear classifier probes'', arXiv, 2016.
\end{omsources}

\section{Exercises}

\begin{exercise}[$\star$]\label{exo:ml:representation-learning:1}
A batch has $m = 256$ examples, two views each. Among how many candidates must each view find its partner under InfoNCE,
and what is the loss if the encoder's codes carry no information?
\end{exercise}

\begin{exercise}[$\star$]\label{exo:ml:representation-learning:2}
How many labelled windows are 300 windows of half a second each, in minutes of trading? And 4\,444?
\end{exercise}

\begin{exercise}[$\star$]\label{exo:ml:representation-learning:3}
An encoder has a 400-dimensional input, 64 hidden units and a 16-dimensional code. How many parameters has it, and how
many has a \omterm{def:ml:representation-learning:probe}{linear probe} on its code?
\end{exercise}

\begin{exercise}[$\star\star$]\label{exo:ml:representation-learning:4}
Why does the \omterm{def:ml:cross-sectional-deep-models:ae}{autoencoder}'s code do worse than the raw window for the forecast, by \cref{prop:ml:represent:pca}?
\end{exercise}

\begin{exercise}[$\star\star$]\label{exo:ml:representation-learning:5}
Why does ridge on the raw window catch up with the predictive probe when the labelled data grow?
\end{exercise}

\begin{exercise}[$\star\star$]\label{exo:ml:representation-learning:6}
\emph{Find the flaw.} ``We pre-trained our encoder on all our data, including the test months, since \omterm{def:ml:representation-learning:ssl}{pre-training} uses no
labels.''
\end{exercise}

\begin{exercise}[$\star\star\star$]\label{exo:ml:representation-learning:7}
\emph{Coding.} Pre-train the predictive encoder to forecast the snapshot's change eight seconds ahead (\texttt{AHEAD = 16})
and rerun \texttt{compare(1000)}. Does the probe improve?
\end{exercise}

\begin{exercise}[$\star\star\star$]\label{exo:ml:representation-learning:8}
Show that the InfoNCE loss of a batch is the cross-entropy of a $(2m - 1)$-way classification and that it equals
$\ln(2m - 1)$ when all similarities are equal.
\end{exercise}

\section{Problem: Learning Without Labels}

\begin{problem}[{Weekend problem --- three pretexts and a few labels}]\label{pb:ml:representation-learning:1}
The chapter's sessions: sixteen unlabelled, four labelled, one for \omterm{def:ml:trees-and-boosting:early}{early stopping}, six for testing.

\medskip\noindent\textbf{Part I --- The pretexts.}
\begin{enumerate}
\item What does each of the three encoders learn to do?
\item Which of them is closest to the task, and why?
\item What does \cref{prop:ml:represent:pca} predict for the \omterm{def:ml:cross-sectional-deep-models:ae}{autoencoder}?
\item What prior do the contrastive augmentations encode?
\end{enumerate}

\medskip\noindent\textbf{Part II --- Few labels.}
\begin{enumerate}[resume]
\item What do the predictive probe, ridge on the raw window and the network from scratch score with 300 labelled windows?
\item And with 1\,000?
\item How do the \omterm{def:ml:cross-sectional-deep-models:ae}{autoencoder} and contrastive probes compare?
\item Why are the hand-made features the worst at 300 and the best at 1\,000?
\end{enumerate}

\medskip\noindent\textbf{Part III --- Many labels.}
\begin{enumerate}[resume]
\item What do the models score with all 4\,444 labelled windows?
\item Why does \omterm{def:ml:representation-learning:ssl}{fine-tuning} help there while the probe does not?
\item What is the fine-tuned model's gain over the network from scratch?
\item What would you expect with forty labelled sessions?
\end{enumerate}

\medskip\noindent\textbf{Part IV --- The verdict.}
\begin{enumerate}[resume]
\item State the \emph{named result}: the labelled sample size below which \omterm{def:ml:representation-learning:ssl}{pre-training} helps, and the probe's IC against the
from-scratch model's at that size.
\item When is \omterm{def:ml:representation-learning:ssl}{pre-training} worth its cost on a desk?
\item How would you choose augmentations for order books?
\item What must be kept out of the \omterm{def:ml:representation-learning:ssl}{pre-training} data?
\item How would you pre-train across instruments?
\item What does a \omterm{def:ml:representation-learning:probe}{linear probe} tell you that \omterm{def:ml:representation-learning:ssl}{fine-tuning} does not?
\item Where does a pre-trained language model fit in this picture (chapter~14)?
\item In one sentence: what makes a pretext useful?
\end{enumerate}
\end{problem}

\section{Interview questions}

\begin{interviewq}[$\star$ \iqroles{mle}]\label{iq:ml:representation-learning:1}
What is the difference between a \omterm{def:ml:representation-learning:probe}{linear probe} and \omterm{def:ml:representation-learning:ssl}{fine-tuning}, and when do you use each?
\end{interviewq}

\begin{interviewq}[$\star\star$ \iqroles{mle, researcher}]\label{iq:ml:representation-learning:2}
Explain \omterm{def:ml:representation-learning:contrastive}{contrastive learning}. What would you use as augmentations for intraday price data?
\end{interviewq}

\begin{interviewq}[$\star\star$ \iqroles{researcher}]\label{iq:ml:representation-learning:3}
You have ten years of unlabelled tick data and three months of labels for a new strategy. How do you use the ten years?
\end{interviewq}

\begin{interviewq}[$\star\star$ \iqroles{researcher, mle}]\label{iq:ml:representation-learning:4}
Why can an \omterm{def:ml:cross-sectional-deep-models:ae}{autoencoder}'s code be useless for prediction although it reconstructs the input well?
\end{interviewq}

\begin{interviewq}[$\star\star$ \iqroles{mle}]\label{iq:ml:representation-learning:5}
How do you evaluate a pre-trained representation without committing to a downstream model?
\end{interviewq}

\begin{interviewq}[$\star\star\star$ \iqroles{researcher}]\label{iq:ml:representation-learning:6}
Prove that the best linear \omterm{def:ml:cross-sectional-deep-models:ae}{autoencoder} with a $k$-dimensional code projects onto the first $k$ principal components.
\end{interviewq}
